\documentclass[12pt]{article}
\usepackage{fontspec}\usepackage{xeCJK}\usepackage{geometry}
\usepackage{fancyhdr}\usepackage{titlesec}\usepackage{caption}
\usepackage{graphicx}\usepackage{amsmath}\usepackage{booktabs}
\usepackage{tikz}\usetikzlibrary{positioning,arrows.meta,fit,calc,backgrounds}
\usepackage[RPvoltages]{circuitikz}
\usepackage{listings}\usepackage{xcolor}

% Rust 代码清单样式（仅 ASCII，避免 CJK 进 listing）
\definecolor{cmt}{RGB}{110,120,130}
\definecolor{kw}{RGB}{20,70,160}
\definecolor{str}{RGB}{160,60,40}
\lstdefinestyle{rust}{
  basicstyle=\ttfamily\footnotesize, keywordstyle=\color{kw}\bfseries,
  commentstyle=\color{cmt}\itshape, stringstyle=\color{str},
  numbers=left, numberstyle=\tiny\color{cmt}, numbersep=6pt,
  showstringspaces=false, breaklines=true, columns=fullflexible,
  frame=single, framesep=4pt, xleftmargin=14pt, aboveskip=6pt, belowskip=2pt,
  morekeywords={let,mut,fn,for,in,if,else,while,loop,struct,impl,pub,const,
    as,return,self,Self,f32,i32,u32,i16,u16,usize,i64,u64}}

\geometry{a4paper, portrait, top=3cm, bottom=2cm, left=2.5cm, right=2.5cm,
  headheight=0pt, headsep=0pt, footskip=1.2cm}

\setmainfont{Times New Roman}
\setCJKmainfont{STSong} % macOS 无 SimSun；STSong 同为宋体。Windows/Overleaf 可改回 SimSun
\xeCJKsetup{AutoFakeBold=2.5}

\newcommand{\sanhao}{\fontsize{16pt}{22pt}\selectfont}
\newcommand{\sihao}{\fontsize{14pt}{22pt}\selectfont}
\newcommand{\xiaosi}{\fontsize{12pt}{22pt}\selectfont}
\newcommand{\wuhao}{\fontsize{10.5pt}{15pt}\selectfont}

\titleformat{\section}{\sanhao\bfseries}{\thesection}{1em}{}
\titleformat{\subsection}{\sihao\bfseries}{\thesubsection}{1em}{}
\titleformat{\subsubsection}{\xiaosi\bfseries}{\thesubsubsection}{1em}{}
\titlespacing{\section}{0pt}{12pt}{6pt}
\titlespacing{\subsection}{0pt}{6pt}{6pt}

\DeclareCaptionFont{wuhao}{\wuhao}
\captionsetup{font=wuhao, labelsep=quad}
\renewcommand{\figurename}{图}
\renewcommand{\tablename}{表}

% 浮动排布：放宽比例限制，避免表格聚集成难看的纯浮动页
\renewcommand{\textfraction}{0.05}
\renewcommand{\topfraction}{0.95}
\renewcommand{\bottomfraction}{0.95}
\renewcommand{\floatpagefraction}{0.85}

\pagestyle{fancy}\fancyhf{}
\renewcommand{\headrulewidth}{0pt}
\fancyfoot[R]{\wuhao\thepage}


\begin{document}
\xiaosi

% ===== 封面（按学院模板：校徽 + 题名 + 设计报告 + 信息栏） =====
\thispagestyle{empty}
\begin{center}
  \vspace*{0.5cm}
  \includegraphics[width=0.82\textwidth]{hdu_logo.png}\\[2.8cm]
  {\fontsize{26pt}{36pt}\selectfont\bfseries 单相功率分析仪（B~题）}\\[0.9cm]
  {\fontsize{26pt}{36pt}\selectfont\bfseries 设\ 计\ 报\ 告}\\[3.6cm]
  {\sihao
  \renewcommand{\arraystretch}{1.9}
  \begin{tabular}{@{}p{2.6cm}@{}c@{}}
    学生姓名 & \underline{\makebox[6.5cm]{刘志\quad 凌海逸\quad 邓奥岩}}\\
    组\hspace{2em}号 & \underline{\makebox[6.5cm]{12}}\\
    实验日期 & \underline{\makebox[6.5cm]{2026/07/17}}\\
  \end{tabular}}
\end{center}
\newpage
\setcounter{page}{1}

% ===== 摘要（≤300字，不编号） =====
\begin{center}
  {\sanhao\bfseries 单相交流功率分析仪}\\[6pt]
  {\sanhao\bfseries 摘\quad 要}
\end{center}
\noindent 本装置基于 TI MSPM0G3507 单片机实现单相交流供电系统的功率分析仪。电压经 QNPT02
电压互感器、电流经电流互感器（原方多匝、匝比可在界面设置）实现强弱电隔离，满足安全规范。
片内双 12\,bit ADC 由同一定时器事件（TIMG6 过零）经事件织物同源触发、DMA 相干采样，每工频
周期 128 点、软件锁频，保证 $u$、$i$ 零相位偏斜。微弱电流信号先经\textbf{片内运算放大器}
可编程增益 $\times 8$ 放大再采样，兼顾小电流分辨率与低功耗（无外接运放）。软件（no\_std Rust
$+$ embassy 异步）以时域 $\overline{ui}$ 求有功功率、$U_{\text{rms}}U_{\text{rms}}$ 求视在功率、
$P/S$ 得功率因数，以 q15 定点基-4 FFT 提取 2\,--\,10 次谐波并计算总谐波畸变率（THD）；经对
市售功率分析仪标定（增益、相位、匝比、时钟）后于 128$\times$64 OLED 分页显示。实测：电压与
电流相对误差 $<1\%$，纯阻负载功率因数达 0.999，电流 THD 与标准表偏差 $<2\%$，非休眠功耗实测
$31\sim36$\,mW（$<50$\,mW），各项指标满足要求。

\vspace{6pt}
\noindent{\xiaosi\bfseries 关键词：}功率分析仪；电流互感器；相干采样；谐波分析；片内运放；MSPM0

\newpage

% ===== 正文 =====
\section{方案比较与论证}
\subsection{总体方案}
装置以\textbf{单颗 MSPM0G3507}（Cortex-M0+，双 12\,bit ADC、两路内置运放 OPA、12\,bit DAC、
事件织物与 DMA）为核心，外接互感器传感前端与 OLED 人机界面，由电池经 LDO 供电（图~\ref{fig:block}）。
不用双处理器而选单片机方案，是因为该器件把同步双 ADC、可编程增益运放和参考 DAC 集成在
一起：$u$、$i$ 的零偏斜同步采样与微弱电流放大都在片内完成，省下了外接运放的静态功耗，
$50$\,mW 指标才有余量。信号链为：市电经电压/电流互感器隔离、调理、
偏置至 $V_{DDA}/2$ 后进 ADC；片内相干采样 $\to$ DMA 成帧 $\to$ 时域 RMS/功率、频域 FFT 谐波
$\to$ 标定 $\to$ OLED 显示。

\begin{figure}[htbp]\centering
\begin{tikzpicture}[
  font=\wuhao,
  blk/.style={draw, rounded corners=2pt, align=center, minimum height=9mm,
    inner sep=3pt, fill=black!3},
  iso/.style={draw, dashed, rounded corners=2pt, align=center, minimum height=9mm,
    inner sep=3pt, fill=blue!4},
  >={Stealth[length=2mm]}, line width=0.4pt]
  \node[blk] (mains) at (0,0) {市电\\220\,V};
  \node[iso] (vt)  at (3.0,0.85)  {电压互感器\\QNPT02$+$RC};
  \node[iso] (ct)  at (3.0,-0.85) {电流互感器\\原方 $n$ 匝};
  \node[blk] (opa) at (5.6,-0.85) {片内 OPA\\PGA $\times 8$};
  \node[blk, minimum height=20mm, text width=27mm] (mcu) at (8.9,0)
    {\textbf{MSPM0G3507}\\事件触发双 ADC\\DMA 相干采样\\时域 $P/PF$、q15 FFT};
  \node[blk] (oled) at (12.1,0) {SSD1306\\OLED 显示};
  \node[blk] (tim)  at (5.6,-2.8)  {TIMG6 过零事件};
  \node[blk] (batt) at (8.3,-2.8)  {电池};
  \node[blk] (ldo)  at (10.5,-2.8) {LDO 3.3\,V};

  \draw[->] (mains.north) |- (vt.west);
  \draw[->] (mains.south) |- (ct.west);
  \draw[->] (ct) -- (opa);
  \draw[->] (vt.east)  -- (mcu.west|-vt.east);
  \draw[->] (opa.east) -- (mcu.west|-opa.east);
  \draw[->] (mcu) -- (oled);
  \draw[->] (tim.north) |- (mcu.south west);
  \draw[->] (batt) -- (ldo);
  \draw[->] (ldo.north) -- (ldo.north|-mcu.south);
\end{tikzpicture}
  \caption{系统框图：互感器隔离前端、片内 OPA/双 ADC、DSP 与显示、电池供电及功耗测量接口}
  \label{fig:block}
\end{figure}

\subsection{传感与隔离方案}
按题目安全要求，电压、电流测量\textbf{必须采用隔离方案}。电压侧选 QNPT02 电流取样型电压
互感器（原方串限流电阻取小电流、副方隔离输出），电流侧选穿孔式电流互感器。互感器天然实现
强弱电隔离，且线性度在本量程（$4$\,A$/250$\,V）内不是精度瓶颈——真正的精度由\emph{标定}保证。

\subsection{电流灵敏度方案比较}
电流量程 $4$\,A、但常测小电流（下探 $0.1$\,A 级），互感器副方信号微弱，是精度难点。\\
方案一（增大取样电阻）：副方电压 $=I_{\text{sec}}\times R$，增大 $R$ 可提幅度，但会抬高铁芯
磁通把互感器\emph{顶入饱和}产生谐波失真，且失真随负载漂移，实测将电压 THD 从 $6.5\%$ 恶化到
$15\%$。\\
方案二（原方多匝缠绕）：被测导线以 $n$ 匝穿过互感器，安匝数提高 $n$ 倍，灵敏度线性提升；匝比
$n$ 于界面可设（$1$\,--\,$10$），软件按 $I=I_{\text{meas}}/n$ 还原，正合题目"原方多匝"之意。\\
方案三（片内 OPA 程控增益）：电流信号经片内运放非反相 PGA 放大 $\times 8$ 再进 ADC，等效
分辨率由约 $14$\,LSB 提升至约 $110$\,LSB，且运放在片内、功耗仅几十 $\mu$A。\\
\textbf{综合选择方案二 $+$ 方案三并用}：多匝提升安匝灵敏度、OPA 提升 ADC 量化分辨率，二者
互补；取样电阻则\emph{取小}以避免方案一的饱和失真。

\subsection{同步采样方案比较}
功率与功率因数对 $u$、$i$ 的\emph{相位对齐}极敏感（$1^\circ$ 相位差即可使 $PF\!\approx\!0.5$ 的
功率误差超 $1\%$）。\\
方案一（单 ADC 分时切换 $u$、$i$）：两通道存在固定采样偏斜，等效相位误差，直接劣化 PF。\\
方案二（双 ADC 同源触发）：ADC0 采 $u$、ADC1 采 $i$，由同一定时器过零事件经事件织物\emph{同时}
触发，零通道切换偏斜。\\
综合选择方案二；器件选型坚持双 ADC，原因即在此。

\subsection{功率与谐波算法方案比较}
有功功率：方案一（频域谱线求和）依赖频率精确锁定，失锁时符号与数值均出错；方案二（时域
$P=\frac{1}{N}\sum u[k]i[k]$）对任意波形精确、恒满足 $|P|\le S$、与锁频无关。选方案二。\\
谐波/THD：方案一（加窗 FFT）通用但有窗泄漏与幅度校正误差；方案二（相干采样 FFT）令每周期
严格 $128$ 点，谐波正好落在整数 FFT 频点上，取模即幅值、无需加窗，前提是\emph{软件锁频}使
采样率跟随工频。选方案二，并以电压基波邻频点能量比插值估计工频、动态微调采样定时器实现锁相。

\section{参数计算、系统设计与实现}
\subsection{理论分析与参数计算}
设一帧含 $N$ 点、去直流后的电压电流序列为 $u[k]$、$i[k]$，则数字测量式为
\begin{equation}
  U_{\text{rms}}=\sqrt{\tfrac1N\textstyle\sum_k u[k]^2},\quad
  I_{\text{rms}}=\sqrt{\tfrac1N\textstyle\sum_k i[k]^2},\quad
  P=\tfrac1N\textstyle\sum_k u[k]\,i[k],
\end{equation}
\begin{equation}
  S=U_{\text{rms}}I_{\text{rms}},\qquad
  PF=\frac{P}{S},\qquad
  Q=\sqrt{S^2-P^2}.
\end{equation}
功率因数取 $P/S$ 定义（而非 $\cos\varphi$），对非正弦负载亦成立。谐波经 512 点实数 FFT 得
各次幅值 $I_h$，总谐波畸变率
\begin{equation}
  \text{THD}=\frac{\sqrt{\sum_{h\ge 2} I_h^{2}}}{I_1}.
\end{equation}
个别谐波显示至 $10$ 次（满足要求）；THD 求和\emph{延伸至 32 次}以对齐市售功率分析仪口径——
截断在 $10$ 次会对充电器等强畸变负载系统性偏低（实测 $43\%$ 对 $47\%$）。

\textbf{相干采样}\quad 令采样率 $f_s=128\,f_{\text{mains}}$，则每帧 $512$ 点恰为 $4$ 个工频
周期、基波落于第 $4$ 号频点、$h$ 次谐波落于第 $4h$ 号。工频漂移时由 FFT 邻频点能量比反推
实际每周期点数，按比例微调触发定时器周期使其收敛到 $128$。

\textbf{互感器工作点}\quad 电压侧 $220$\,V 经原方 $110$\,k$\Omega$ 限流得 $I_{\text{pri}}=2$\,mA
（额定值），副方以\textbf{小取样电阻/跨阻}取出、使铁芯磁通远离饱和区（这是电压 THD 达标的
关键，见 \S3.3）。电流侧原方 $n$ 匝，软件量程 $=4$\,A$/n$。

\textbf{电流通道 OPA 与偏置}\quad 片内 OPA1 配为非反相 PGA（$\times 8$），反馈梯形网络底部
经 \texttt{msel} 接片内 DAC12（置 $V_{DDA}/2$），故增益作用于 $(V_{in}-V_{DDA}/2)$，即围绕
中轨放大交流分量而不使直流削顶；开机以 DAC 伺服将输出直流拉至中点，自动吸收前端偏置失配。

\subsection{硬件电路设计}
\textbf{电压通道}：QNPT02 原方串 $110$\,k$\Omega$ 接 $220$\,V；副方取样后经 $R=10$\,k$\Omega$、
$C=1.5$\,nF 抗混叠 RC（截止 $\approx 10.6$\,kHz，对 $10$ 次谐波以内近乎透明）与 $V_{DDA}/2$
偏置进 PA25（ADC0 通道 2）。\\
\textbf{电流通道}：电流互感器原方 $n$ 匝，副方取样电阻取小以避饱和，偏置 $V_{DDA}/2$ 后接
PB19（OPA1\_IN0+）；OPA1 以 $\times 8$ PGA、DAC12 作参考，输出经 ADC1 内部通道 13 采样。\\
\textbf{时序与采集}：TIMG6 以 $32$\,MHz 下计数产生过零事件，经事件织物两通道\emph{同时}触发
ADC0、ADC1，各由一路 DMA 搬入 $512$ 点缓冲；ADC 开 $16\times$ 硬件平均降噪。\\
\textbf{显示与供电}：SSD1306 $128\times 64$ OLED（$\mathrm{I^2C}$，PB2/PB3）分页显示，两枚按键
翻页与设置匝比；$18650$ 电池经 TPS7A4701 低噪声 LDO 得 $3.3$\,V（附录 A.3），电池回路留
mA/V 表 $4$\,mm 香蕉插座测功耗；测评时拔除 LaunchPad 调试器跳线以隔离其功耗。

\subsection{软件设计}
固件为 no\_std Rust $+$ embassy 异步，分四模块。\textbf{采样引擎}：配置事件触发双 ADC 与
DMA、逐帧锁频（调 TIMG6 周期）。\textbf{DSP}：时域 RMS/有功功率、q15 基-4 512 点 FFT（板载
实测约 $4.8$\,ms/帧）提取谐波与 THD、邻频点插值估工频；含孤立尖峰中值去毛刺与 8 帧箱型滑动
平均。\textbf{标定}：结构 \texttt{Calib}$=\{k_u,k_i,\varphi,\text{sign},f_{clk}\}$，对市售
PA 标定电压/电流标度、传感相位失配、CT 极性与时钟偏差。\textbf{界面}：OLED 分页（功率/电网/
谐波/柱状图/设置），按键设置匝比 $n$。为规避片上 $\mathrm{I^2C}$ 驱动缺陷，OLED 采用每事务
$\le 8$ 字节的分块写。系统以采集协程与显示协程经信号量交互（图~\ref{fig:ui}）。

\begin{figure}[htbp]\centering
{\setlength{\fboxsep}{0pt}\fbox{\includegraphics[width=0.46\textwidth]{oled_power.png}}}
  \caption{OLED 主页（Power 页）：以固件渲染代码在宿主机逐像素输出的 $128\times64$
    位图（最近邻放大），内容为表~\ref{tab:acc} 的纯阻负载读数}
  \label{fig:ui}
\end{figure}

\section{测量方案与测量结果分析}
\subsection{测试仪器}
市售单相功率分析仪（作标定与比对基准）、数字示波器 Siglent SDS2012X Plus（观测互感器波形、
定位失真）、调压器、直流 mA/V 表；测试负载取白炽灯（纯阻）、白炽灯并联电容构成的容性负载，
以及开关电源充电器（强畸变负载，仅作趋势验证）。

\subsection{测试方案与数据}
\textbf{标定}：以 PA 读数为真值，分别标定电压标度 $k_u$（电压有效值比）、电流标度 $k_i$（电流
有效值比，与匝比解耦）、传感链相位失配 $\varphi$、CT 极性与时钟 $f_{clk}$，系数固化于固件。
\textbf{比对}：白炽灯纯阻负载给出稳定基准（表~\ref{tab:acc}，谐波见表~\ref{tab:harm}）；
白炽灯并联电容的容性负载（$PF\approx0.87$）检验非单位功率因数下的功率分解与谐波测量
（表~\ref{tab:cap}，谐波见表~\ref{tab:capharm}）；充电器为非平稳负载仅作趋势验证。

\begin{table}[htbp]\centering
\caption{纯阻负载（白炽灯）下各参数与市售功率分析仪（PA）比对}
\label{tab:acc}
\wuhao
\begin{tabular}{lccc}
\toprule
参数 & PA 基准 & 本装置 & 相对误差\\
\midrule
电压 $U$\,(V)      & 223.1  & 223.0 & $<1\%$\\
电流 $I$\,(A)      & 0.273  & 0.273 & $<1\%$\\
有功功率 $P$\,(W)  & 61.0   & 60.7  & $-0.5\%$\\
功率因数 $PF$      & 0.999  & 0.999 & ---\\
电流 THD          & 6.0\%  & 6.4\% & $+0.4\%$（绝对）\\
\bottomrule
\end{tabular}
\end{table}

\begin{table}[htbp]\centering
\caption{电流逐次谐波幅值（mA）与 PA 对照（纯阻负载，白炽灯）}
\label{tab:harm}
\wuhao
\begin{tabular}{lcccccccccc}
\toprule
谐波次数 & 1 & 2 & 3 & 4 & 5 & 6 & 7 & 8 & 9 & 10\\
\midrule
PA\,(mA)    & 272.5 & 0.3 & 11.7 & 0 & 3.3 & 0 & 6.3 & 0   & 4.4 & 0\\
本装置\,(mA) & 272.4 & 1.1 & 14.7 & 0 & 1.9 & 0 & 4.9 & 0.3 & 5.2 & 0\\
误差\,(mA)   & $-0.1$ & $+0.8$ & $+3.0$ & 0 & $-1.4$ & 0 & $-1.4$ & $+0.3$ & $+0.8$ & 0\\
\bottomrule
\end{tabular}

\vspace{2pt}
\end{table}

\begin{table}[htbp]\centering
\caption{容性负载（白炽灯并联电容）下各参数与 PA 比对}
\label{tab:cap}
\wuhao
\begin{tabular}{lccc}
\toprule
参数 & PA 基准 & 本装置 & 相对误差\\
\midrule
电压 $U$\,(V)      & 218.6  & 218.53 & $-0.03\%$\\
电流 $I$\,(A)      & 0.310  & 0.306  & $-1.3\%$\\
有功功率 $P$\,(W)  & 58.9   & 58.5   & $-0.7\%$\\
功率因数 $PF$      & 0.868  & 0.872  & $+0.004$（绝对）\\
电流 THD          & 20.8\% & 20.1\% & $-0.7\%$（绝对）\\
\bottomrule
\end{tabular}
\end{table}

\begin{table}[htbp]\centering
\caption{电流逐次谐波幅值（mA）与 PA 对照（容性负载）}
\label{tab:capharm}
\wuhao
\begin{tabular}{lcccccccccc}
\toprule
谐波次数 & 1 & 2 & 3 & 4 & 5 & 6 & 7 & 8 & 9 & 10\\
\midrule
PA\,(mA)    & 302 & 0   & 24   & 0 & 16   & 0 & 19   & 0   & 20   & 0\\
本装置\,(mA) & 299 & 0.8 & 22.2 & 0 & 18.8 & 0 & 16.8 & 0.1 & 19.6 & 0.4\\
误差\,(mA)   & $-3$ & $+0.8$ & $-1.8$ & 0 & $+2.8$ & 0 & $-2.2$ & $+0.1$ & $-0.4$ & $+0.4$\\
\bottomrule
\end{tabular}

\vspace{2pt}
{\wuhao 注：本组按 PA 的幅值显示口径记录；PA 谐波分辨率 $1$\,mA，有效位比本装置少一位。}
\end{table}

\subsection{结果分析与结论}
表~\ref{tab:acc} 表明电压、电流有效值相对误差均 $<1\%$（要求 1），有功功率与功率因数误差
$<1\%$（要求 2），电流 THD 与 PA 偏差 $0.4\%$（$<2\%$，要求 3）。表~\ref{tab:harm} 表明各次
谐波幅值与 PA 逐次吻合，偶次基本为零、奇次量级一致；电流与电压谐波\emph{并不完全相同}（3 次谐波
电流占比 $5.4\%$，即 $14.7/272.4$\,mA，对电压 $6.1\%$），恰说明本装置测得的是电流自身谐波而非电压耦合，与 PA "电流 THD 略低于
电压" 的方向一致。容性负载（表~\ref{tab:cap}、表~\ref{tab:capharm}）进一步检验了非单位
功率因数工作点：$PF$ 偏差仅 $0.004$（绝对），验证了相位标定在 $PF\ne 1$ 下依然成立；THD
$20.1\%$ 对 $20.8\%$、逐次谐波（3/5/7/9 次为主）幅值一致；本装置同时给出 $S=66.2$\,VA、
$Q=32.2$\,var，与 $S=UI$、$Q=\sqrt{S^2-P^2}$ 自洽。

本设计有三处关键工程要点：其一，\textbf{电流灵敏度}靠"原方多匝 $+$ 片内 OPA 增益"解决，使
$0.27$\,A 小电流的信号从约 $14$\,LSB 提升到约 $110$\,LSB，小电流下 PF/谐波方达标；其二，
\textbf{电压 THD} 曾达 $15\%$，经示波器定位为电压互感器副方取样电阻过大致铁芯饱和（偶次谐波
显著），\emph{减小取样电阻}使其退出饱和后 THD 回落至电网真实值 $\sim6.5\%$——软件不可逆地
修复硬非线性，只能改硬件工作点；其三，\textbf{相位标定}须在信号充分（大信号）的纯阻负载上做，
且以频域相量旋转而非时域插值，方不损失相关性。功耗方面，电流放大用\emph{片内}运放（几十
$\mu$A）而非外接运放，主控降频运行（$32$\,MHz，不启用 $80$\,MHz PLL，省去 PLL 自身电流），实测电池端（$4.11$\,V）电流 $7.54\sim8.76$\,mA、
非休眠总功耗 $31\sim36$\,mW（$<50$\,mW，要求 4）。综上，本装置在市电
单相供电下完成电流、电压、有功功率、功率因数、THD 及 2\,--\,10 次谐波的测量与显示，全部指标
满足题目要求。

\appendix
\section{附录 A\quad 前端电路原理图}
\subsection*{A.1\quad 电压测量通道}
电压互感器原方经 $110$\,k$\Omega$ 限流取 $2$\,mA（额定），副方以 $100$\,$\Omega$ 小取样电阻
取出（避免铁芯饱和），经 $10$\,k$\Omega$/$1.5$\,nF 抗混叠 RC 滤波后，在 ADC 节点以
$10$\,k$\Omega$（上拉 $3.3$\,V）/$100$\,k$\Omega$（下拉）注入偏置，进 ADC0（图~\ref{fig:vch}）。
偏置网络与信号回路的直流内阻（$\approx 10$\,k$\Omega$）合成，节点静态电位约在半轨。

\begin{figure}[htbp]\centering
\begin{circuitikz}[font=\wuhao, american, scale=0.92, transform shape]
  \node[transformer] (T) at (3.0,0) {};
  \node[above=1mm of T] {QNPT02};
  \draw (0,0.8) to[sV, l=$220$\,V] (0,-0.8);
  \draw (0,0.8) to[R, l=$110\,$k$\Omega$] ++(1.9,0) -| (T.A1);
  \draw (0,-0.8) -| (T.A2);
  % 副方 100R 取样
  \coordinate (bt) at ($(T.B1)+(1.0,0)$);
  \coordinate (bb) at (bt|-T.B2);
  \draw (T.B1) -- (bt) to[R, l_=$100\,\Omega$, *-*] (bb);
  \draw (T.B2) -- (bb);
  % RC 抗混叠
  \coordinate (c) at ($(bt)+(2.4,0)$);
  \draw (bt) to[R, l=$10\,$k$\Omega$, -*] (c);
  \draw (c) to[C, l_=$1.5\,$nF] (c|-bb);
  \draw (bb) -- (c|-bb) node[ground]{};
  % 偏置网络（ADC 节点）
  \coordinate (o) at ($(c)+(1.8,0)$);
  \draw (c) to[short, -*] (o);
  \draw (o) to[R, l_=$10\,$k$\Omega$] ++(0,1.7) node[above]{$3.3$\,V};
  \draw (o) to[R, l=$100\,$k$\Omega$] ++(0,-1.7) node[ground]{};
  \draw (o) to[short, -o] ++(1.3,0) node[right]{PA25 (ADC0)};
\end{circuitikz}
  \caption{电压测量通道：电压互感器 $+$ 原方限流 $+$ 副方 $100\,\Omega$ 取样 $+$
    抗混叠 RC $+$ 偏置网络}
  \label{fig:vch}
\end{figure}

\subsection*{A.2\quad 电流测量通道}
电流互感器原方 $n$ 匝，副方 $100$\,$\Omega$ 取样、$10$\,k$\Omega$/$1.5$\,nF 抗混叠 RC 与
同规格偏置网络（$10$\,k$\Omega$ 上拉/$100$\,k$\Omega$ 下拉）后送 PB19，进片内运放 OPA1
非反相 PGA（$\times 8$，反馈梯形网络底部以片内 DAC12 作 $V_{DDA}/2$ 参考），输出经 ADC1
内部通道 13 采样（图~\ref{fig:ich}）。

\begin{figure}[htbp]\centering
\begin{circuitikz}[font=\wuhao, american, scale=0.92, transform shape]
  \node[transformer] (CT) at (2,0) {};
  \node[above=1mm of CT] {CT};
  \draw (0,0.8) node[left]{被测线（$n$ 匝）} -| (CT.A1);
  \draw (0,-0.8) -| (CT.A2);
  % 副方 100R 取样
  \coordinate (bt) at ($(CT.B1)+(1.0,0)$);
  \coordinate (bb) at (bt|-CT.B2);
  \draw (CT.B1) -- (bt) to[R, l_=$100\,\Omega$, *-*] (bb);
  \draw (CT.B2) -- (bb);
  % RC 抗混叠
  \coordinate (c) at ($(bt)+(2.4,0)$);
  \draw (bt) to[R, l=$10\,$k$\Omega$, -*] (c);
  \draw (c) to[C, l_=$1.5\,$nF] (c|-bb);
  \draw (bb) -- (c|-bb) node[ground]{};
  % 偏置网络
  \coordinate (p) at ($(c)+(1.8,0)$);
  \draw (c) to[short, -*] (p);
  \draw (p) to[R, l_=$10\,$k$\Omega$] ++(0,1.7) node[above]{$3.3$\,V};
  \draw (p) to[R, l=$100\,$k$\Omega$] ++(0,-1.7) node[ground]{};
  % 片内 OPA（同相输入对齐信号线）
  \node[op amp, scale=0.85] (OA) at ($(p)+(2.8,0.42)$) {};
  \draw (p) -- ($(p)+(1.3,0)$) node[above]{PB19} -| (OA.+);
  \draw (OA.out) to[short, -o] ++(1.0,0) node[right]{ADC1 ch13};
  \node at ($(OA.center)+(1mm,-1pt)$) {$\times8$};
  \node[below=9mm of OA, align=center] {DAC12 $=V_{DDA}/2$（参考）};
  \begin{scope}[on background layer]
    \node[draw, dashed, gray, rounded corners, fit=(OA), inner sep=6pt,
      label={[gray]above:片内}] {};
  \end{scope}
\end{circuitikz}
  \caption{电流测量通道：电流互感器（原方 $n$ 匝）$+$ 副方 $100\,\Omega$ 取样 $+$
    抗混叠 RC $+$ 偏置网络 $+$ 片内 OPA PGA（$\times8$，DAC 参考）}
  \label{fig:ich}
\end{figure}

\subsection*{A.3\quad 供电与功耗测量}
$18650$ 锂电池（$3.7\sim4.2$\,V）经 TPS7A4701 低噪声 LDO 稳压至 $3.3$\,V。输出电压由
接地设定引脚编程：$V_{\text{OUT}}=1.4+1.6+0.2+0.1=3.3$\,V（跳线接地 1P6V/0P2V/0P1V，
SENSE 经 $0$\,$\Omega$ 直连输出）。电池回路引出 $4$\,mm 香蕉插座，串入电流表、并联电压表
直接测整机功耗；测评时拔除 LaunchPad 调试器跳线以隔离 XDS110（图~\ref{fig:pwr}）。

\begin{figure}[htbp]\centering
\begin{circuitikz}[font=\wuhao, american, scale=0.92, transform shape]
  % 电池 + 功耗测量回路
  \draw (0,0) to[battery1, invert] (0,-2.4) node[ground]{};
  \node[left, align=center] at (-0.35,-1.2) {18650\\ $3.7\sim4.2$\,V};
  \draw (0,0) to[ammeter] (2.2,0);
  \node[above] at (1.1,0.55) {功耗测量（香蕉插座）};
  \draw (2.2,0) to[short, -*] (3.2,0);
  \draw (3.2,0) to[voltmeter] (3.2,-2.4) node[ground]{};
  % 输入电容
  \draw (3.2,0) to[short, -*] (4.6,0);
  \draw (4.6,0) to[C, l=$10\,\mu$F] (4.6,-1.5) node[ground]{};
  % LDO 框（IN 引脚对齐 y=0；加高留足引脚行距）
  \node[draw, minimum width=4.0cm, minimum height=3.4cm] (U) at (8.4,-1.2) {};
  \node[above, font=\wuhao\bfseries] at (U.north) {TPS7A4701};
  \coordinate (in)  at ($(U.west)+(0,1.2)$);
  \coordinate (en)  at ($(U.west)+(0,0.1)$);
  \coordinate (nr)  at ($(U.west)+(0,-1.0)$);
  \coordinate (out) at ($(U.east)+(0,1.2)$);
  \coordinate (fb)  at ($(U.east)+(0,0.1)$);
  \coordinate (set) at ($(U.south)+(-1.0,0)$);
  \coordinate (gnd) at ($(U.south)+(1.4,0)$);
  \node[right, font=\wuhao] at (in) {IN};
  \node[right, font=\wuhao] at (en) {EN};
  \node[right, font=\wuhao] at (nr) {NR/SS};
  \node[left, font=\wuhao]  at (out) {OUT};
  \node[left, font=\wuhao]  at (fb) {SENSE};
  \node[above, font=\wuhao] at ($(set)+(0,0.05)$) {1P6V\,0P2V\,0P1V};
  \node[above, font=\wuhao] at ($(gnd)+(0,0.05)$) {GND};
  % 输入侧连线：IN 直连；EN 上拉至 IN 网络
  \draw (4.6,0) -- (in);
  \draw (en) -- ++(-0.6,0) -- (5.8,0|-en) (5.8,0|-en) -- (5.8,0) node[circ]{};
  % NR/SS 软启动电容（外移，避开底部设定引脚文字）
  \draw (nr) -- ++(-0.9,0) to[C, l_=$1\,\mu$F] ++(0,-0.9) node[ground]{};
  % 输出侧：SENSE 0R 直连 OUT，两颗 22uF
  \draw (out) -- (11.6,0);
  \draw (fb) -- ++(0.6,0) -- (11.0,0|-fb) (11.0,0|-fb) -- (11.0,0) node[circ]{};
  \draw (11.6,0) to[C, l=$22\,\mu$F$\times 2$, *-] (11.6,-1.5) node[ground]{};
  \draw (11.6,0) to[short, -o] (12.7,0) node[right]{$3.3$\,V 系统};
  % 设定引脚（接地编程 3.3V）与地
  \draw (set) -- ++(0,-0.45) node[ground]{};
  \node[below, font=\wuhao] at ($(set)+(0,-0.8)$) {跳线接地};
  \draw (gnd) -- ++(0,-0.45) node[ground]{};
\end{circuitikz}
  \caption{供电与功耗测量：$18650$ 电池 $+$ 串联电流表/并联电压表（香蕉插座）$+$
    TPS7A4701 LDO（设定引脚接地编程 $3.3$\,V）}
  \label{fig:pwr}
\end{figure}

\section{附录 B\quad 关键源码（Rust，no\_std）}
固件（Rust，\texttt{no\_std}$+$embassy 异步）按功能分四个模块：采样（事件触发双 ADC 与
DMA 相干采样、锁频）、测量（时域 RMS/功率、q15 基-4 FFT 谐波/THD、工频估计与标定）、界面
（OLED 显示与按键）、主控（OPA/DAC 配置、DC 伺服与主循环）。核心测量与标定代码摘录如下。

\begin{lstlisting}[style=rust,caption={测量核心摘录：时域 RMS/有功功率/功率因数与相干 FFT 的 THD}]
// One frame (N=512, 4 mains cycles) -> U/I RMS, active power, PF, THD.
let u_mean = mean(&u_raw);
let i_mean = mean(&i_raw);
let (mut su2, mut si2, mut sui) = (0i64, 0i64, 0i64);
for k in 0..N {
    let du = (u_raw[k] & 0xFFF) as i32 - u_mean; // remove DC bias
    let di = (i_raw[k] & 0xFFF) as i32 - i_mean;
    su2 += (du * du) as i64;
    si2 += (di * di) as i64;
    sui += (du * di) as i64;
}
let urms = sqrt(su2 as f32 / N as f32) * cal.ku;
let irms = sqrt(si2 as f32 / N as f32) * cal.ki / turns;
// Active power in the time domain: exact for any waveform, |P| <= S.
let p = (sui as f32 / N as f32) * cal.ku * (cal.ki / turns) * cal.isign;
let s = urms * irms;
let pf = (p / s).clamp(-1.0, 1.0);

// Harmonics via coherent q15 radix-4 512-pt FFT: h-th harmonic on bin 4h.
let spec = fft512_q15(&centered_i);
let fund = mag(&spec, 4);
let mut acc = 0.0;
for h in 2..=32 {
    acc += mag(&spec, 4 * h) * mag(&spec, 4 * h); // THD summed to the 32nd
}
let thd = sqrt(acc) / fund; // display H2..H10 separately
\end{lstlisting}

\begin{lstlisting}[style=rust,caption={标定系数与相干采样/锁频摘录}]
/// Calibrated once against a reference power analyzer, then frozen in
/// firmware.
pub struct Calib {
    pub ku: f32,    // V per ADC LSB (voltage scale)
    pub ki: f32,    // A per ADC LSB (current, before /turns)
    pub phase: f32, // u -> i sensor phase mismatch (rad)
    pub isign: f32, // CT polarity (+1 / -1)
    pub fclk: f32,  // sample-clock (SYSOSC) correction
}

/// 128 pts/cycle, 4 cycles/frame (N = 512).
const SAMPLES_PER_CYCLE: usize = 128;

// TIMG6 zero-event -> event fabric ch1/ch2 -> ADC0(u) & ADC1(i) at the
// SAME instant (zero skew). Each frame, retune the timer so pts/cycle
// converges to 128:
fn track(&mut self, spc: f32) {
    let period = (self.load + 1) as f32;
    let target = period * spc / SAMPLES_PER_CYCLE as f32;
    self.load = (period + 0.6 * (target - period)) as u32 - 1; // damped lock
}
\end{lstlisting}

\end{document}
